\documentclass[10pt]{article}
\providecommand{\FIGDIR}{fig}
\newcommand{\DIGESTDATE}{17 September 2026}
% ===================== PM Research Watch - shared preamble =====================
\usepackage[T1]{fontenc}
\usepackage[utf8]{inputenc}
\usepackage{amsmath,amssymb}
\usepackage{tgpagella}
\usepackage{tgheros}
\renewcommand{\sfdefault}{qhv}
\renewcommand{\ttdefault}{lmtt}
\usepackage{microtype}
\usepackage[english]{babel}

\usepackage[a4paper,top=17mm,bottom=17mm,left=17mm,right=17mm,
            headsep=6mm,footskip=10mm]{geometry}
\usepackage{graphicx}
\usepackage{xcolor}
\usepackage{colortbl}
\usepackage{tikz}
\usetikzlibrary{calc}
\usepackage[most]{tcolorbox}
\usepackage{titlesec}
\usepackage{fancyhdr}
\usepackage{booktabs}
\usepackage{array}
\usepackage{longtable}
\usepackage{enumitem}
\usepackage{multicol}
\usepackage{ragged2e}
\usepackage[hidelinks]{hyperref}

% ------------------------------------------------------------------ palette
\definecolor{Ink}    {HTML}{14262E}
\definecolor{Deep}   {HTML}{0F4C5C}
\definecolor{Teal}   {HTML}{1F7A8C}
\definecolor{Sky}    {HTML}{4FA3C4}
\definecolor{Sage}   {HTML}{6FA287}
\definecolor{Amber}  {HTML}{E3A008}
\definecolor{Clay}   {HTML}{B5643C}
\definecolor{Coral}  {HTML}{D1495B}
\definecolor{Violet} {HTML}{6A5B9E}
\definecolor{Slate}  {HTML}{4A6572}
\definecolor{Paper}  {HTML}{FBF9F5}
\definecolor{Mist}   {HTML}{EEEAE2}
\definecolor{Rule}   {HTML}{D9D3C9}

\pagecolor{Paper}
\color{Ink}

\hypersetup{colorlinks=true, linkcolor=Deep, urlcolor=Teal, citecolor=Teal,
            pdftitle={PM Research Watch}, pdfauthor={Automated daily digest}}

% ------------------------------------------------------------------ headings
\titleformat{\section}
  {\sffamily\bfseries\large\color{Deep}}{}{0pt}
  {}[{\vspace{1.2mm}\color{Amber}\titlerule[1.1pt]}]
\titlespacing*{\section}{0pt}{5.5mm}{2.8mm}

\titleformat{\subsection}
  {\sffamily\bfseries\normalsize\color{Teal}}{}{0pt}{}
\titlespacing*{\subsection}{0pt}{3.6mm}{1.4mm}

\setlist[itemize]{leftmargin=4.2mm, itemsep=1.1mm, topsep=1.3mm, parsep=0pt,
                  label=\textcolor{Teal}{\footnotesize$\blacktriangleright$}}

% ------------------------------------------------------------------ header/footer
\pagestyle{fancy}
\fancyhf{}
\renewcommand{\headrulewidth}{0pt}
\renewcommand{\footrulewidth}{0pt}
\fancyhead[L]{\sffamily\scriptsize\color{Slate}PM RESEARCH WATCH\;\textcolor{Rule}{\textbar}\;\DIGESTDATE}
\fancyhead[R]{\sffamily\scriptsize\color{Slate}Particulate matter monitoring \& health effects}
\fancyfoot[C]{\sffamily\scriptsize\color{Slate}\thepage}

% ------------------------------------------------------------------ components
% metric tile
\newcommand{\metric}[3]{%
  \begin{tcolorbox}[enhanced, colback=white, colframe=Rule, boxrule=0.5pt,
    arc=1.6mm, left=2.2mm, right=2.2mm, top=1.8mm, bottom=1.8mm,
    borderline west={1.7mm}{0pt}{#3}, nobeforeafter, valign=center,
    height=17mm, width=\linewidth]
    \RaggedRight
    {\sffamily\bfseries\fontsize{17}{17}\selectfont\color{#3}#1}\\[0.5mm]
    {\sffamily\fontsize{6.5}{7.4}\selectfont\color{Slate}#2}
  \end{tcolorbox}}

% pull-quote / key insight box
\newtcolorbox{keybox}[1][]{enhanced, colback=Mist, colframe=Mist,
  arc=1.6mm, left=3.2mm, right=3.2mm, top=2.4mm, bottom=2.4mm,
  borderline west={1.7mm}{0pt}{Amber}, #1}

% subtopic band
\newcommand{\band}[2]{%
  \vspace{3.2mm}\noindent
  \begin{tcolorbox}[enhanced, colback=#2, colframe=#2, arc=1.2mm,
    left=2.6mm, right=2.6mm, top=1.1mm, bottom=1.1mm, boxrule=0pt,
    width=\linewidth, nobeforeafter]
    {\sffamily\bfseries\small\color{white}#1}
  \end{tcolorbox}\vspace{1.6mm}}

% one paper entry: {short cite}{title}{journal}{doi}{body}
\newcommand{\paperentry}[5]{%
  \par\vspace{1.9mm}\noindent
  {\sffamily\bfseries\footnotesize\color{Ink}#1}\;%
  {\sffamily\fontsize{7}{8}\selectfont\color{Slate}#3}\\[0.4mm]
  {\footnotesize\itshape\color{Deep}#2}\\[0.6mm]
  {\footnotesize #4}\\[0.3mm]
  {\fontsize{6.6}{7.6}\selectfont\color{Slate}\texttt{doi:}\,\href{https://doi.org/#5}{\texttt{\detokenize{#5}}}}
  \par}

% caption style
\newcommand{\figcap}[1]{\par\vspace{1.0mm}
  {\fontsize{7.2}{8.6}\selectfont\color{Slate}\textsf{\textbf{Fig.}}\; #1\par}\vspace{1.5mm}}

\providecommand{\permil}{\textperthousand}

\setlength{\parindent}{0pt}
\setlength{\parskip}{0pt}
\renewcommand{\arraystretch}{1.18}


\begin{document}

% ==================================================================== MASTHEAD
\thispagestyle{empty}
\begin{tikzpicture}[remember picture, overlay]
  \fill[Deep] (current page.north west) rectangle
      ($(current page.north east)+(0,-67mm)$);
  \fill[Amber] ($(current page.north west)+(0,-67mm)$) rectangle
      ($(current page.north east)+(0,-68.4mm)$);
  \foreach \x/\y/\r in {22/12/0.9, 41/26/1.5, 63/9/0.7, 88/31/1.1, 112/17/1.9,
                        138/29/0.8, 159/13/1.3, 178/54/1.0, 196/21/1.6,
                        30/44/0.6, 74/57/1.2, 125/48/0.9, 168/6/0.8, 191/59/0.7,
                        52/18/1.0, 100/60/0.6, 148/20/1.1, 66/49/1.4, 205/45/0.8}
    \fill[Sky, opacity=0.30] ($(current page.north west)+(\x mm,-\y mm)$) circle (\r mm);
\end{tikzpicture}

\vspace*{4mm}
{\sffamily\fontsize{7.6}{9}\selectfont\color{Sky}%
 AUTOMATED DAILY LITERATURE DIGEST\;\textbullet\;PhD RESEARCH WATCH\par}
\vspace{2.2mm}
{\sffamily\bfseries\fontsize{27}{29}\selectfont\color{white}%
 Particulate Matter\par}
{\sffamily\fontsize{27}{29}\selectfont\color{Sky}%
 Monitoring \& Health Effects\par}
\vspace{5.0mm}
{\sffamily\fontsize{8.4}{10}\selectfont\color{white}%
 Issue 2026-09-17\;\;\textcolor{Amber}{\textbar}\;\;PubMed entry window 17 September 2026\;\;%
 \textcolor{Amber}{\textbar}\;\;18 records screened in scope\par}
\vspace{18mm}

% -------------------------------------------------------------- metric strip
\noindent
\begin{minipage}[t]{0.190\linewidth}\metric{18}{RECORDS IN SCOPE\\(57 EXCLUDED, 5 HELD)}{Deep}\end{minipage}\hfill
\begin{minipage}[t]{0.190\linewidth}\metric{7}{OF 10 SUBTOPIC\\BANDS POPULATED}{Teal}\end{minipage}\hfill
\begin{minipage}[t]{0.190\linewidth}\metric{4}{EFFECT ESTIMATES\\(ALL WITH A CI)}{Sky}\end{minipage}\hfill
\begin{minipage}[t]{0.190\linewidth}\metric{6}{TIER-A\\RECORDS}{Amber}\end{minipage}\hfill
\begin{minipage}[t]{0.190\linewidth}\metric{0}{CONSENSUS RECORDS\\ADMITTED}{Clay}\end{minipage}

\vspace{6mm}

% -------------------------------------------------------------- signal of day
\section*{\sffamily\bfseries\color{Deep}Signal of the day}
\vspace{-1.5mm}{\color{Amber}\rule{\linewidth}{1.1pt}}\vspace{2.5mm}

\begin{keybox}
{\sffamily\bfseries\small\color{Deep}A 115{,}056-birth null: short-term PM$_{2.5}$ shows
no association with spontaneous preterm birth in two independent US cohorts, under three
lag specifications, with intervals tight enough that the null is informative rather than
merely underpowered.}\\[1.4mm]
\footnotesize
Nelin et al.\ (\textit{Ann Epidemiol}) analysed \textbf{109{,}735 spontaneous preterm
births in Michigan (2003--2020)} and \textbf{5{,}321 in metropolitan Philadelphia
(2008--2020)} with a time-stratified case-crossover design against daily high-resolution
modelled PM$_{2.5}$ in the week before delivery. They ran the exposure three ways ---
nested daily lag terms, averaged lag exposures, and distributed lag non-linear models ---
and used likelihood ratio tests to check whether the flexible lag structure improved fit.
\textbf{No specification, in either cohort, produced a statistically significant
association}, with point estimates near the null and, in the authors' words, narrow
confidence intervals. The precision is what makes this worth reading: a null with wide
intervals says nothing, and a two-cohort null replicated across three lag
parameterisations is a much harder result to dismiss than a single flat coefficient.
Two caveats bound it. First, exposure is \emph{modelled} at the residential address, not
measured, and non-differential exposure error biases precisely toward this answer --- the
paper does not quantify that error, which is the one thing that would settle whether the
null is biological or methodological. Second, case-crossover self-matching removes
time-invariant confounding but nothing time-varying beyond the temperature adjustment
applied. The authors are explicit that the result holds \textbf{at the studied
concentrations} and press for continued surveillance as regulatory protections change.
\textbf{For instrumentation the reading is direct:} if modelled exposure error is what
separates this null from the positive short-term literature, then the binding constraint
on reproductive air-pollution epidemiology is measurement resolution at the address, not
another cohort.
\end{keybox}

\clearpage

% -------------------------------------------------------------- exec brief
\section*{\sffamily\bfseries\color{Deep}What changed today}
\vspace{-1.5mm}{\color{Amber}\rule{\linewidth}{1.1pt}}\vspace{2.5mm}

\footnotesize
\begin{itemize}
\item \textbf{Three leading US energy-system models disagree by a factor of 2.6 on the same
decarbonisation health question.} The \textit{ES\&T} multi-model assessment (US-REGEN,
GCAM-USA, REPEAT) finds PM$_{2.5}$-attributable mortality rates in 2050 fall
\textbf{28--73\%} against 2017 under net-zero. Direction is robust; magnitude is not, and
the spread tracks how each model allocates mitigation across sectors --- \textbf{carbon
dioxide removal deployment in 2050 varies threefold}, setting how much residual fossil
combustion and precursor emission survives. The sensitivity analysis is the sharper
finding: plausible increases in non-energy emissions, wildfire in particular, can
\textbf{offset the energy-sector benefit outright}. Note what the 28--73\% band does
\emph{not} contain --- all three models share concentration-response functions, so the
published spread understates total uncertainty.

\item \textbf{Clean air and thermal comfort are close to anti-correlated across four Asian
residential sites.} The \textit{AAQR} sensor-driven framework crossed a sustained
PM$_{2.5}$ screen (trailing 24-h mean, 15\,$\mu$g/m$^3$ WHO threshold) with a 24--30\,\textdegree C
/ RH\,$\le$\,65\% envelope over 724--744 hourly bins per site. Clean-hour fractions ran
\textbf{14.5 / 34.3 / 30.8 / 97.4\%} for Kuala Lumpur, Bogor, Singapore and Palmerston,
while standalone thermal feasibility ran \textbf{97.0 / 71.5 / 98.5 / 20.6\%} --- the site
with clean air has almost no comfortable hours, and the reverse. A natural-ventilation
policy set on either constraint alone inverts at three of the four sites. The data are
uncalibrated AirGradient nodes read through OpenAQ with no co-location against a reference
monitor, which is the obvious weakness and also the realistic deployment case.

\item \textbf{A consumer sleep app is now a population respiratory endpoint.} Xia et al.\
(\textit{Commun Health}) ran an AI cough detector over nocturnal audio from Sleep Cycle
users across \textbf{32 cities in 12 countries over ${\sim}$500 days}: each 10\,$\mu$g/m$^3$
PM$_{2.5}$ gave a \textbf{2.4\% rise in nocturnal cough frequency (RR 1.024, 95\% CI
1.013--1.035)}, non-linear, with measurable increases at low concentrations and a visible
excursion during a wildfire episode. The cough detector is never validated against a
clinical reference in this paper, so outcome misclassification is unquantified --- but as
a demonstration that a passive consumer sensor stream can be paired with a PM network at
city-day resolution it stands on its own.

\item \textbf{Dropping the strongest pollutants made the residual mixture look worse.}
The Beijing COPDB panel (\textit{J Hazard Mater}, 112 older adults, 410 visits) used BKMR
to identify SO$_2$ and UFPs for fasting glucose, PM$_{2.5}$ Ca$^{2+}$ for heart rate, and
O$_3$ and UFPs for reduced peak flow --- then removed them and re-fitted. Residual
associations \emph{strengthened}: the gaseous group over 14 days tied to
\textbf{+19.1\% fasting glucose (10.8--27.4)} and \textbf{$-$16.0\% peak flow (8.7--23.3)}.
Read as a warning about single-pollutant screening it is fair; read as a finding it is
partly a collinearity artefact of the removal step, and $n=112$ is thin for pairwise
interaction across a dozen-plus exposures.

\item \textbf{Consensus returned nothing admissible for the second consecutive issue ---
this time across two distinct axes.} Twenty records over a calibration sweep and a
satellite-fusion sweep: \textbf{six were already in the archive} by title, and
\textbf{all fourteen survivors failed a Crossref \texttt{created} check}, the earliest
2026-01-02 and the latest 2026-07-06. Nothing came within ten weeks of the window. Two
axes failing identically moves this from a one-off to a property of the source.

\item \textbf{Five tier-A-plausible records were held, not shipped.} All five are
same-day deposits with no abstract anywhere --- three Elsevier (\textit{Atmos Environ}
WRF-Chem PM$_{2.5}$ forecasting via deep-learning emission inversion; \textit{STOTEN} Doha
compound heat-pollution; \textit{Atmos Res} physics-based aerosol source apportionment at
Lecce) and two Frontiers/Elsevier records with no DOI deposited at all. This is the
\textbf{third consecutive issue} losing on-axis candidates to the no-abstract-on-deposit
pattern. All five are kept out of \texttt{seen.json} so a later issue can carry them.
\end{itemize}

\clearpage

% ==================================================================== ANALYTICS
\section*{\sffamily\bfseries\color{Deep}Corpus analytics}
\vspace{-1.5mm}{\color{Amber}\rule{\linewidth}{1.1pt}}\vspace{3mm}

\noindent\includegraphics[width=\linewidth]{\FIGDIR/f1_subtopics.png}
\figcap{Subtopic assignment of the 18 in-scope records; single-label by dominant
contribution. \textit{Exposure assessment \& modelling} dominates with \textbf{6}, a third
of the issue and twice any other band. \textit{Sensing, forecasting \& instrumentation} and
\textit{Burden, policy \& mitigation} hold 3 each, \textit{Respiratory \& allergic} and
\textit{Neuro / mental health} 2 each, and \textit{Reproductive \& developmental} and
\textit{Cardiovascular \& metabolic} 1 each. Together the two measurement-side bands hold
\textbf{9 of 18 --- exactly half the issue}. \textbf{Three bands are empty}:
\textit{Mechanistic toxicology}, \textit{Occupational \& indoor} and \textit{Other clinical
endpoints}, so 7 of 10 are populated.}

\vspace{2mm}
\noindent\includegraphics[width=\linewidth]{\FIGDIR/f2_design_endpoint.png}
\figcap{Left --- study architecture. \textit{Measurement campaign} is the largest group at
\textbf{6 of 18 (33\%)}, with \textit{Modelling / inventory} and
\textit{Observational --- acute} at 3 each, \textit{Review / synthesis} at 2, and
\textit{Experimental / toxicology}, \textit{Observational --- cohort},
\textit{Metadata only} and \textit{Ecological} at 1 each. There is \textbf{no randomised
trial and no cross-sectional study} in this issue. Right --- health endpoint.
\textit{No health endpoint} is the largest single bar at \textbf{5}, followed by
\textit{Respiratory} at 3 and \textit{Instrument performance} at 2; \textbf{eight further
endpoint classes hold exactly one record each}. Counting the instrument-performance and
air-quality bars alongside the explicit no-endpoint one, \textbf{8 of 18 records carry no
human health outcome at all}.}

\vspace{2mm}
\noindent\includegraphics[width=\linewidth]{\FIGDIR/f3_metric_geography.png}
\figcap{Left --- particle metric. \textit{PM$_{2.5}$ only} leads at \textbf{9}, exactly
half the issue, with \textit{PM (unspeciated) / emissions} at 7 and \textit{optical
properties} and \textit{PM$_{2.5}$ + PM$_{10}$ jointly} at 1 each --- a narrower metric
spread than recent issues, with no bioaerosol or composition-resolved record. Right ---
geographic provenance. \textit{Global / multi-region} leads with \textbf{6} and
\textit{China} follows with \textbf{5}, so the two together carry 11 of 18;
\textit{North America} holds 3, and Southeast Asia, Sub-Saharan Africa, South Asia and
Europe hold one each. \textbf{Latin America and the Middle East \& N. Africa are absent.}}

\vspace{2mm}
\noindent\includegraphics[width=\linewidth]{\FIGDIR/f6_lifecourse.png}
\figcap{Life-course windows addressed; counts are non-exclusive. The
\textbf{preconception / prenatal window leads with 3} --- the 115{,}056-birth
case-crossover, the Colorado gestational PM$_{2.5}$--birth-weight application, and the
prenatal arm of the corrected \textit{BBI-H} study. The childhood and adolescent windows
each show 1, but \textbf{both rest entirely on that same corrigendum}, which is
metadata-only: there is no primary paediatric or adolescent paper in this issue, and the
count should not be read as though there were. Older adults hold 2 (the 3{,}617-participant
Hubei cognition cohort and the 112-adult Beijing panel). \textbf{The working-age window is
empty}, consistent with the \textit{Occupational \& indoor} band being unpopulated.}

\vspace{2mm}
\noindent\includegraphics[width=\linewidth]{\FIGDIR/f4_heatmap.png}
\figcap{Subtopic $\times$ study-architecture cross-tabulation. Three features are worth
naming precisely. First, the \textit{Measurement campaign} column is
\textbf{entirely measurement-side}: all 6 of its records sit in \textit{Sensing} (2) or
\textit{Exposure assessment} (4), with none on any health band. Second, of the 6
\textit{Exposure assessment} records \textbf{exactly one carries a human health endpoint}
--- the Colorado multi-tree birth-weight application; the other five are aerosol physics,
aerosol chemistry, source apportionment, formation chemistry and an ozone-concentration
analysis. Third, \textit{Observational --- acute} is spread \textbf{one record per band}
across \textit{Respiratory}, \textit{Reproductive} and \textit{Cardiovascular}, so no
health band is carried by more than a single acute design. \textit{Mechanistic toxicology}
is absent entirely this issue, so the usual all-experimental row does not appear.}

\vspace{2mm}
\noindent\includegraphics[width=\linewidth]{\FIGDIR/f5_forest.png}
\figcap{The 4 interval-bearing estimates in this issue, on a common logarithmic axis.
\textbf{None of the four crosses unity.} \textbf{Only one is a PM exposure contrast} ---
nocturnal cough per 10\,$\mu$g/m$^3$ PM$_{2.5}$, RR 1.024 (1.013--1.035) --- reversing the
15 September issue, whose forest plot was carried by PM, and returning to the pattern that
has held for most of the archive. The other three are not PM at all: the cognitive-impairment
OR 2.18 (1.31--3.64) is per unit of a \emph{compound heat} index (Humidex$_{\mathrm{mean}}$)
from a study in which PM$_{2.5}$ enters only as a covariate, and the two Beijing estimates
--- fasting glucose 1.191 (1.108--1.274) and peak expiratory flow 0.840 (0.767--0.913) ---
are \textbf{gaseous-mixture} contrasts over a 14-day window, converted here from the
published percent changes to ratio form and flagged as such. \textbf{Exposure contrasts are
therefore not mutually comparable}: read direction and precision, not relative magnitude.
The cough estimate's tightness rests on ${\sim}$500 city-days across 32 cities; the Beijing
intervals rest on 410 visits from 112 people and should be read as the least stable of the
four despite not being the widest.}

\clearpage

% ==================================================================== DIGEST
\section*{\sffamily\bfseries\color{Deep}Paper-level digest}
\vspace{-1.5mm}{\color{Amber}\rule{\linewidth}{1.1pt}}\vspace{1mm}

{\fontsize{7.4}{9}\selectfont\color{Slate}Ordered by cluster, largest first. Each entry
gives the methodological core, the finding that matters and the limitation that actually
threatens the inference --- not an abstract paraphrase.\par}

\band{Neuro / mental health\hfill 2 records}{Deep}

\paperentry{Hubei Memory \& Aging 2026}
{Associations between compound heat exposure and cognitive impairment and the potential mediating role of plasma sphingolipid metabolism.}
{Environ Pollut}
{3,617 non-dementia adults aged 65+ in the Hubei Memory and Aging Cohort, repeatedly assessed 2016-2025. Excess cumulative temperature in heatwaves was linearly related to cognitive impairment for daytime-only, nighttime-only and compound heatwaves; temperature-humidity indices gave near-U-shaped curves with Humidexmean best fitting (OR 2.18, 95\% CI 1.31-3.64). Among 77 plasma sphingolipid measures, only the ceramide C22:0 / sphingomyelin C26:0 ratio mediated, and at -14.3\% - a suppression effect, not an explanatory pathway, which the abstract soft-pedals. PM$_{2.5}$ enters only as a covariate, so this is a heat paper that adjusts for particles rather than a PM paper. Testing 77 metabolites for a single mediation signal invites a multiplicity problem that is not addressed.}
{10.1016/j.envpol.2026.129182}

\paperentry{Corrigendum 2026}
{Corrigendum to 'Prenatal and childhood air pollution exposure, cellular immune biomarkers, and brain connectivity in early adolescence'.}
{Brain Behav Immun Health}
{\textsc{metadata only.} Corrigendum notice only; the record carries no abstract beyond the pointer to the 2024 article it corrects, and the deposited text does not state which results changed. Logged rather than summarised because the corrected paper sits squarely on a line this digest tracks - modelled early-life PM$_{2.5}$ against immune biomarkers and adolescent brain connectivity - and because a correction to that specific inference chain is worth flagging to anyone citing the 2024 original. No effect estimate can be extracted from the notice, and the direction and magnitude of the correction remain unknown until the corrected version is read directly. Treat any citation of the 2024 estimates as provisional until then.}
{10.1016/j.bbih.2026.101317}

\band{Cardiovascular \& metabolic\hfill 1 record}{Teal}

\paperentry{COPDB Beijing 2026}
{Associations between air pollution and cardiorespiratory function and metabolic indicators: New insights into mixture exposures.}
{J Hazard Mater}
{112 older adults in the Beijing COPDB study across 410 follow-up visits. Single-pollutant BKMR flagged SO$_2$ and UFPs for fasting blood glucose, PM$_{2.5}$ Ca2+ for heart rate, and O$_3$ and UFPs for reduced peak expiratory flow. The novel step is removing those pollutants and re-fitting: the residual mixture associations got stronger, with the gaseous group over a 14-day window tied to a 19.1\% FBG increase (10.8-27.4\%) and a 16.0\% PEF decrease (8.7-23.3\%). That is a real warning about single-pollutant screening, but the design invites it - dropping the strongest predictors and finding the rest matter more is partly a collinearity artefact, and n=112 with 410 visits is thin for pairwise interactions across a dozen-plus exposures. Component-resolved measurement is the strength.}
{10.1016/j.jhazmat.2026.143550}

\band{Respiratory \& allergic\hfill 2 records}{Sky}

\paperentry{Xia et al. 2026}
{Short-term associations between air pollution and nocturnal cough frequency.}
{Commun Health}
{Nightly cough counts inferred by an AI cough detector from Sleep Cycle audio, aggregated to city-day across 32 cities in 12 countries over \textasciitilde{}500 days. Each 10 $\mu$g\,m$^{-3}$ PM$_{2.5}$ was associated with a 2.4\% rise in nocturnal cough frequency (RR 1.024, 95\% CI 1.013-1.035), with a non-linear exposure-response showing increases at low concentrations and a visible excursion during a wildfire episode. Two threats dominate: the cough detector's accuracy is never validated against a clinical reference in this paper, so outcome misclassification is unquantified, and Sleep Cycle users are a self-selected smartphone-owning population whose composition varies by city. Relevant as proof that a consumer sensor stream can act as a population respiratory endpoint alongside a PM network.}
{10.1038/s44528-026-00030-5}

\paperentry{Korean J Intern Med 2026}
{Health effects of air pollution on chronic respiratory diseases.}
{Korean J Intern Med}
{Narrative synthesis of the epidemiological evidence tying air pollution to incidence, exacerbation and mortality across asthma, COPD, interstitial lung disease and lung cancer, then a section on non-pharmacological exposure-reduction measures at individual and societal level. No search strategy, no inclusion criteria, no quantitative pooling, so it cannot be read as evidence in its own right - the value is the intervention section, which is the part least covered by the recent primary literature. Included as a reference point for what the field currently treats as settled rather than for any new estimate. The interventions discussed (purifiers, mask use, activity rescheduling) all presuppose exposure information at a spatial resolution regulatory networks do not provide.}
{10.3904/kjim.2025.423}

\band{Reproductive \& developmental\hfill 1 record}{Sage}

\paperentry{Nelin et al. 2026}
{Associations of short-term changes in air pollution with spontaneous preterm birth: A time-stratified case-crossover analysis.}
{Ann Epidemiol}
{A well-powered null: 109,735 spontaneous preterm births in Michigan (2003-2020) and 5,321 in metropolitan Philadelphia (2008-2020), analysed with nested daily lags, averaged lags and DLNMs, all temperature-adjusted, with likelihood ratio tests comparing linear-lag and DLNM fits. No significant association emerged in either cohort under any specification, with point estimates near null and narrow intervals - the narrowness is what makes this informative rather than merely underpowered. The exposure is modelled, not measured, so non-differential error at the residential address would bias exactly toward this result, and the authors do not quantify that error. Case-crossover also removes only time-invariant confounding. Worth pairing with the positive short-term literature rather than reading alone.}
{10.1016/j.annepidem.2026.110291}

\band{Exposure assessment \& modelling\hfill 6 records}{Clay}

\paperentry{Multi-Tree 2026}
{Multi-Tree Model for Precision Environmental Health with Longitudinally Assessed Mixture Exposure.}
{J R Stat Soc Ser C Appl Stat}
{A Bayesian ensemble of tree triplets - two interacting binary tree types - that estimates individualised and subgroup distributed-lag effects for mixture exposures assessed repeatedly, with built-in modifier selection over a high-dimensional candidate set. Applied to Colorado birth registry data it recovers a PM$_{2.5}$-birth weight association concentrated in mothers under 25 with low income. The applied finding is the weaker half: effect-modification by maternal age and income is already well established, so it functions as a sanity check rather than a discovery, and no interval is reported for the subgroup estimate in the abstract. The methodological contribution is the part worth using - most distributed-lag mixture methods return one population-average curve, and this returns a per-individual one.}
{10.1093/jrsssc/qlag047}

\paperentry{Sichuan Basin 2026}
{Topographic Modulation of High-Temperature Ozone Suppression: Core-Periphery Threshold Structure in the Sichuan Basin.}
{Environ Pollut}
{High-temperature suppression - ozone flattening or falling under extreme heat rather than rising with the usual climate penalty - appeared in 16 of 18 analysed Sichuan Basin cities. The critical suppression temperature fell from 36.8 C in the basin interior to 33.8 C at the periphery, and the decline rate above it reached -9.65 $\mu$g\,m$^{-3}$/C in the inner zone, about 2.1x the outer. SHAP attributions put boundary-layer ventilation ahead of cloud and circulation terms. Ozone, not PM, so it sits at the edge of scope; included because the same basin ventilation structure governs PM accumulation and because the threshold-by-topography result is a caution against fitting a single regional temperature response. Five years of data and a closed basin limit transfer elsewhere.}
{10.1016/j.envpol.2026.129161}

\paperentry{ORACLES 4STAR 2026}
{Subseasonal and spatial variability of biomass burning aerosol radiative properties observed over the Southeast Atlantic during ORACLES.}
{Atmos Chem Phys}
{All three ORACLES deployment months (2016-2018) of 4STAR sky-scans over the Southeast Atlantic smoke-over-stratocumulus regime. Single scattering albedo rises through the burning season - the plume brightens - and the absorption Angstrom exponent analysis gives no evidence of a brown carbon contribution, so the authors attribute it to composition change rather than a change of aerosol type. Comparison against a 30-year climatology from 31 southern African AERONET sites shows the same seasonal direction over land and sea. Sampling is the caveat the paper handles well: the September SSA dip is shown to be an artefact of 2016's more southerly flight tracks. Three campaign years cannot separate a seasonal cycle from interannual variability, which they concede.}
{10.5194/acp-26-13083-2026}

\paperentry{Delhi N2O5 2026}
{Enhanced nocturnal N2O5-ClNO$_2$ chemistry in Delhi.}
{Atmos Chem Phys}
{Early-spring 2023 measurements in Delhi against a winter 2019 baseline. Nocturnal N2O5 and ClNO$_2$ averaged 13 and 80 ppt, roughly double 2019, and the authors attribute this to nighttime NO falling from 124+/-25 to 44+/-9 ppb, which removes the main N2O5 sink. Particulate chloride at 4.7 $\mu$g\,m$^{-3}$ is among the highest reported anywhere and drives efficient N2O5-to-ClNO$_2$ conversion, raising NO$_3$ and Cl radical production and organic oxidation. The comparison is the weak joint - different seasons (spring vs winter) and different years are confounded, so the NO-decline explanation competes with a straightforward seasonal one. Matters for PM because this chemistry feeds nitrate and secondary organic mass, the fraction low-cost optical sensors systematically mishandle under high humidity.}
{10.5194/acp-26-13121-2026}

\paperentry{China NOx isotopes 2026}
{Driving factors of atmospheric nitrogen oxides emissions and pollution in China.}
{Environ Pollut}
{Twenty-two years of nitrate isotope data put fossil-fuel combustion at \textasciitilde{}53\% of national NOx and non-fossil sources (biomass burning, microbial N cycling) at 47\%. The paper is unusually candid about what it cannot resolve: regional source-fraction differences 'lack statistical significance', and the isotope-derived NOx trend is a non-significant 7\% decline (p=0.19), which the authors say precludes robust inference of emission-deposition decoupling - even though a 35\% post-2011 drop in surface NO$_2$ against only 19\% less nitrate deposition is the headline. A NOy deposition saturation threshold of 15.0 kg N/ha/yr is proposed. Relevant to PM as nitrate is a major secondary component; the saturation claim rests on a fitted inflection with no uncertainty stated.}
{10.1016/j.envpol.2026.129167}

\paperentry{PRD limona ketone 2026}
{Ground-gridded observations of high concentrations of ambient limona ketone suggest overlooked limonene oxidation in subtropical megacities.}
{Environ Pollut}
{Large-scale gridded plus intensive diurnal measurements of limona ketone across the Pearl River Delta found unusually high concentrations during spring and autumn ozone episodes. A parameterised inversion using the product as a constraint implies more limonene precursor than inventories carry, with a pronounced nighttime enhancement; modelled production is OH-dominated while the loss pathway shifts from OH by day to the nitrate radical at night. The inference is indirect by construction - precursor abundance is inverted from a product, so it inherits every error in the assumed yield and loss rates, and the authors only claim the nocturnal enhancement is robust 'across plausible parameter ranges' rather than quantifying it. Episode-specific sampling limits it further. Matters as biogenic SOA is an under-constrained part of urban PM mass.}
{10.1016/j.envpol.2026.129177}

\band{Sensing, forecasting \& instrumentation\hfill 3 records}{Amber}

\paperentry{Zhang et al. 2026}
{Adaptive robotic sensing for exposure-oriented mapping of hazardous airborne particle dispersion in indoor environments.}
{J Hazard Mater}
{CFD virtual sampling plus a full-scale mechanically ventilated aerosol release. The claim worth taking is the region-of-interest one: adaptive sampling cut normalised reconstruction error in the high-concentration ROI to 9.5\% at 65 points against 18.3\% for uniform scanning, and cut mean error at four fixed reference locations from 59.7\% to 34.3\% in the physical experiment. Note the honest counter-result the authors report - optimised uniform scanning beat adaptive sampling on final global NRMSE (2.5\% vs 3.0\%), so the gain is purely in plume localisation, not in field-wide accuracy. Single ventilation condition, single room, no repeat releases, so the 34.3\% figure has no dispersion around it. Directly relevant to mobile-node deployment: it quantifies what a moving sensor buys over a fixed sparse array for exactly the transient indoor plume case fixed networks miss.}
{10.1016/j.jhazmat.2026.143612}

\paperentry{Front Bioeng 2026}
{Bionic breathing-type airway exposure system for inhalation exposure studies.}
{Front Bioeng Biotechnol}
{A 3D-printed trachea-to-G3 airway replica with millisecond breathing control (five waveforms, oral/nasal/combined routes), CO$_2$ exchange and multi-zone exposure interfaces. Validation recovers the expected right-dominant airflow split (right main bronchus 51.6+/-0.1\% vs left 48.4+/-0.1\%) and separates deposition by breathing pattern (large-loop 78.0+/-1.6\% vs small-loop 61.7+/-1.3\%). BEAS-2B cells showed dose-dependent cytotoxicity; hiPSC-derived lung models gave selective inflammatory responses. The weakness is that airflow-split fidelity is a low bar - it follows from the geometry, not the dynamics - and no comparison against in vivo human deposition data is offered, so anatomical realism is asserted rather than demonstrated. For sensing work it matters as a dosimetry bench: it converts a measured mass concentration into a branch-resolved deposited dose.}
{10.3389/fbioe.2026.1871869}

\paperentry{AAQR 2026}
{A Sensor-Driven State-Space Framework for PM$_{2.5}$-Informed Natural Ventilation Opportunities Under Coupled Pollution-Thermal-Humidity Constraints.}
{Aerosol Air Qual Res}
{Four residential communities - Kuala Lumpur, Bogor, Singapore and Palmerston (AU) - over 1 Jan to 1 Feb 2026, 724-744 hourly bins each. Crossing a sustained-PM$_{2.5}$ screen with a T-RH envelope (24-30 C, RH <=65\%) shows the two constraints are close to anti-correlated: clean-hour fractions were 14.5 / 34.3 / 30.8 / 97.4\%, while standalone thermal feasibility ran 97.0 / 71.5 / 98.5 / 20.6\%. The site with clean air has almost no comfortable hours and vice versa. Limits are real - one month, one season, four sites, uncalibrated AirGradient data taken as-is with no co-location against reference monitors, and fixed rather than locally derived thresholds. Still the most directly usable result of the issue for LCS network design: it turns a raw sensor feed into an actionable control signal.}
{10.1007/s44408-026-00167-6}

\band{Burden, policy \& mitigation\hfill 3 records}{Coral}

\paperentry{US net-zero multi-model 2026}
{A Multi-Model Assessment of the Air Quality and Health Impacts of U.S. Energy Decarbonization.}
{Environ Sci Technol}
{Three independent US energy system models (US-REGEN, GCAM-USA, REPEAT) run to 2050 agree on direction but not magnitude: PM$_{2.5}$-attributable mortality rates fall 28-73\% against 2017 under net-zero, despite population growth and ageing. The spread is the result. It tracks how each model allocates CO$_2$ mitigation across sectors, with carbon dioxide removal deployment in 2050 varying threefold and setting how much residual fossil combustion - and therefore precursor emission - remains. A sensitivity analysis shows plausible increases in non-energy emissions, wildfire in particular, can offset the energy-sector benefit outright. Concentration-response functions are shared across the three models, so the reported spread understates total uncertainty: model structure varies, epidemiological input does not.}
{10.1021/acs.est.6c02914}

\paperentry{EOTBL cancer 2026}
{Global Burden and Temporal Trends of Early-Onset Tracheal, Bronchial, and Lung Cancer: An Analysis of GLOBOCAN 2022 and GBD 2023 Data.}
{Adv Sci}
{Early-onset tracheal, bronchial and lung cancer caused 137,705 incident cases and 72,646 deaths globally in 2022, with the highest incidence in Eastern Asia and in High-HDI settings, but disproportionate mortality in Medium- and Low-HDI regions - the familiar detection-versus-survival split. The PM-relevant claim is that ambient particulate pollution emerges as a leading attributable risk in rapidly developing economies such as China, where active smoking dominates in countries like Russia. That attribution is GBD machinery, not new observation: it inherits GBD's exposure surfaces and its concentration-response assumptions wholesale, and no uncertainty interval accompanies the risk-factor ranking here. Registry coverage is also weakest in exactly the Low-HDI regions carrying the mortality burden.}
{10.1002/advs.77778}

\paperentry{Campania review 2026}
{Air pollution and chronic respiratory diseases: a narrative review including a case study in Campania (Italy).}
{Monaldi Arch Chest Dis}
{Narrative review of mechanism (oxidative stress, chronic inflammation, epithelial damage, airway remodelling) attached to a Campania case study, one of the Italian regions most persistently above national average. Between 2015 and 2024 national PM$_{10}$ and PM$_{2.5}$ fell slightly while Campania stayed high, with measurable respiratory impact and marked intra-regional disparity that the authors tie to socioeconomic gradient. Closing sections argue for machine learning and satellite-based forecasting to improve monitoring, and for environmental justice targeting. No systematic search, no pooled estimate, and the regional numbers are cited rather than re-analysed. Included for the inequality framing and because the monitoring-gap argument it makes is the case for dense low-cost networks in exactly such regions.}
{10.4081/monaldi.2026.4216}

\clearpage

\section*{\sffamily\bfseries\color{Deep}Record register}
\vspace{-1.5mm}{\color{Amber}\rule{\linewidth}{1.1pt}}\vspace{2.5mm}

{\fontsize{6.7}{8.4}\selectfont
\setlength{\tabcolsep}{2.4pt}
\begin{longtable}{@{}>{\RaggedRight}p{25mm}>{\RaggedRight}p{41mm}>{\RaggedRight}p{27mm}%
>{\RaggedRight}p{23mm}>{\RaggedRight}p{18mm}>{\RaggedRight}p{15mm}@{}}
\toprule
\textsf{\textbf{First author}} & \textsf{\textbf{Focus}} & \textsf{\textbf{Journal}} &
\textsf{\textbf{Design}} & \textsf{\textbf{Metric}} & \textsf{\textbf{Region}}\\
\midrule
\endfirsthead
\toprule
\textsf{\textbf{First author}} & \textsf{\textbf{Focus}} & \textsf{\textbf{Journal}} &
\textsf{\textbf{Design}} & \textsf{\textbf{Metric}} & \textsf{\textbf{Region}}\\
\midrule
\endhead
\multicolumn{6}{@{}l@{}}{\cellcolor{Mist}\textsf{\textbf{\textcolor{Deep}{Neuro / mental health}}}}\\[0.6mm]
\href{https://doi.org/10.1001/jamanetworkopen.2026.25015}{Dales} & Wildfire-sourced PM$_{2.5}$ and ED visits for migraine and other primary headache & \textit{JAMA Netw Open} & Case-crossover & PM$_{2.5}$ (wildfire) & Canada\\
\href{https://doi.org/10.1093/ageing/afag215}{Glover} & Umbrella meta-review of urban environmental exposures and cognitive health & \textit{Age Ageing} & Umbrella review / meta-analysis & PM$_{2.5}$ & Global\\
\href{https://doi.org/10.1016/j.nsa.2026.107019}{Goossens} & In utero air pollution exposure and pre-/postnatal brain development & \textit{Neurosci Appl} & Narrative review & PM$_{2.5}$ / PM$_{10}$ / PM$_{\mathrm{coarse}}$ & Global\\
\href{https://doi.org/10.1016/j.isci.2026.116617}{Ma} & PM$_{2.5}$-induced neurotoxicity and depression: exposure risks and mechanisms & \textit{iScience} & Narrative review & PM$_{2.5}$ + components & Global\\
\href{https://doi.org/10.3760/cma.j.cn112150-20250804-00752}{Pan} & Maternal PM exposure and offspring neurodevelopment: effects and mechanisms & \textit{Zhonghua Yu Fang Yi Xue Za Zhi} & Narrative review & PM (bulk) & China\\
\addlinespace[0.9mm]
\multicolumn{6}{@{}l@{}}{\cellcolor{Mist}\textsf{\textbf{\textcolor{Teal}{Cardiovascular \& metabolic}}}}\\[0.6mm]
\href{https://doi.org/10.3760/cma.j.cn112338-20251120-00843}{Huang} & PM$_{2.5}$ and dyslipidaemia in rural Chinese children; dietary effect modification & \textit{Zhonghua Liu Xing Bing Xue Za Zhi} & Cross-sectional (multistage) & PM$_{2.5}$ & China\\
\href{https://doi.org/10.32394/pe/218966}{Kazmierska} & Air pollution and cardiovascular disease: review of current data & \textit{Przegl Epidemiol} & Narrative review & PM$_{2.5}$ / PM$_{10}$ & Global\\
\href{https://doi.org/10.1007/s00125-026-06814-2}{Xi} & Air pollution, built environment and temperature vs. prediabetes phenotypes & \textit{Diabetologia} & Prospective cohort & PM$_{2.5}$ & Germany\\
\href{https://doi.org/10.1161/JAHA.125.046917}{Yang} & Short-term PM$_{2.5}$ and hospitalisation/ED use in atrial fibrillation & \textit{J Am Heart Assoc} & Case-crossover & PM$_{2.5}$ (wildfire-influenced) & USA\\
\addlinespace[0.9mm]
\multicolumn{6}{@{}l@{}}{\cellcolor{Mist}\textsf{\textbf{\textcolor{Sky}{Reproductive \& developmental}}}}\\[0.6mm]
\href{https://doi.org/10.1016/j.jes.2025.10.043}{Chen} & Prenatal air pollution, fetal growth trajectory phenotypes and child neurodevelopment & \textit{J Environ Sci (China)} & Prospective cohort & PM$_{2.5}$ / PM$_{10}$ & China\\
\href{https://doi.org/10.1016/j.ecoenv.2026.120566}{Huang} & Joint prenatal bisphenol + air-pollutant mixture exposure and preschool IQ & \textit{Ecotoxicol Environ Saf} & Prospective cohort (mixtures) & PM$_{2.5}$ / PM$_{10}$ & China\\
\href{https://doi.org/10.1007/s40572-026-00552-8}{Peterson} & Ambient air pollution, PAHs and pubertal development: critical review & \textit{Curr Environ Health Rep} & Critical review & PM$_{2.5}$ components & Global\\
\href{https://doi.org/10.1016/j.envres.2026.125313}{Xie} & Long-term PM$_{2.5}$ and timing of pubertal development in boys and girls & \textit{Environ Res} & Prospective cohort & PM$_{2.5}$ & China\\
\addlinespace[0.9mm]
\multicolumn{6}{@{}l@{}}{\cellcolor{Mist}\textsf{\textbf{\textcolor{Sage}{Respiratory \& allergic}}}}\\[0.6mm]
\href{https://doi.org/10.1007/s00484-026-03289-z}{Bayar Muluk} & Climate change and otolaryngology: airway disease impacts and emission reduction & \textit{Int J Biometeorol} & Narrative review & PM$_{2.5}$ & Global\\
\href{https://doi.org/10.4103/mgr.MEDGASRES-D-25-00164}{Liao} & Short-term air pollution, respiratory symptoms and lung function in asthma & \textit{Med Gas Res} & Panel study & PM$_{2.5}$ / PM$_{10}$ & China\\
\href{https://doi.org/10.3390/epidemiologia7040097}{Sanchez de Prada} & Ambient PM and asthma-coded admissions in children 0-3 y, Spain & \textit{Epidemiologia} & Ecological panel & PM$_{2.5}$ / PM$_{10}$ & Spain\\
\href{https://doi.org/10.1016/j.jhazmat.2026.143076}{Yang} & Allergenic vs. toxic risk of airborne seed hairs compared with ambient PM$_{2.5}$ & \textit{J Hazard Mater} & Field measurement & PM$_{2.5}$ + bioaerosol & China\\
\addlinespace[0.9mm]
\multicolumn{6}{@{}l@{}}{\cellcolor{Mist}\textsf{\textbf{\textcolor{Violet}{Mechanistic toxicology}}}}\\[0.6mm]
\href{https://doi.org/10.1016/j.jes.2025.12.074}{Ding} & PM$_{2.5}$ promotes vascular calcification via miR-27b-5p/HES1 signalling & \textit{J Environ Sci (China)} & In vitro + organoid + murine & PM$_{2.5}$ & China\\
\href{https://doi.org/10.1007/s00360-026-01698-5}{Maraschi} & Osmoregulatory disruption and metal bioaccumulation from settleable atmospheric PM & \textit{J Comp Physiol B} & Controlled exposure (ecotox) & Settleable APM & Brazil\\
\href{https://doi.org/10.1016/j.jes.2025.11.018}{Wu} & Macrophage NLRP3 inflammasome in real-ambient PM-induced hepatic glucose dysregulation & \textit{J Environ Sci (China)} & Real-ambient murine + co-culture & Real-ambient PM & China\\
\addlinespace[0.9mm]
\multicolumn{6}{@{}l@{}}{\cellcolor{Mist}\textsf{\textbf{\textcolor{Clay}{Occupational \& indoor}}}}\\[0.6mm]
\href{https://doi.org/10.3389/fmedt.2026.1796828}{Chen} & Nested suction instrument to cut intraoperative surgical-smoke PM$_{2.5}$ & \textit{Front Med Technol} & Bench / device experiment & PM$_{2.5}$ (surgical smoke) & China\\
\href{https://doi.org/10.1093/annweh/wxag055}{Cole} & Determinants of respirable crystalline silica exposure in Australian tunnelling & \textit{Ann Work Expo Health} & Observational exposure study & Respirable dust / RCS & Australia\\
\href{https://doi.org/10.1016/j.jhazmat.2026.143081}{Ding} & Stage-resolved hair Zn-isotope dynamics as a firefighting combustion-exposure tracer & \textit{J Hazard Mater} & Repeated-measures biomarker & Combustion-derived PM & China\\
\href{https://doi.org/10.3390/jox16040119}{Ramsden} & Xenobiotic hazards in aircraft cabin air & \textit{J Xenobiot} & Review + risk economics & Ultrafine particles & Global\\
\addlinespace[0.9mm]
\multicolumn{6}{@{}l@{}}{\cellcolor{Mist}\textsf{\textbf{\textcolor{Amber}{Sensing, forecasting \& instrumentation}}}}\\[0.6mm]
\href{https://doi.org/10.1016/j.jenvman.2026.130515}{Shi} & Adaptive spatiotemporal graph transformer for multi-site multi-step PM$_{2.5}$ forecasting & \textit{J Environ Manage} & Deep-learning model & PM$_{2.5}$ & China / India\\
\href{https://doi.org/10.3390/bios16070380}{Tsuchiya} & Seasonal and indoor-outdoor characteristics of fluorescent aerosol particles in offices & \textit{Biosensors} & Real-time sensor campaign & Bioaerosol / coarse & Japan\\
\addlinespace[0.9mm]
\multicolumn{6}{@{}l@{}}{\cellcolor{Mist}\textsf{\textbf{\textcolor{Amber}{Exposure assessment \& modelling}}}}\\[0.6mm]
\href{https://doi.org/10.1029/2025GH001621}{Solorzano-Araque} & Multiscale exposure modelling: LOTOS-EUROS + mobility surveys + particle toxicity & \textit{GeoHealth} & CTM + mobility integration & PM$_{2.5}$ / PM$_{10}$ & Colombia\\
\addlinespace[0.9mm]
\multicolumn{6}{@{}l@{}}{\cellcolor{Mist}\textsf{\textbf{\textcolor{Coral}{Burden, policy \& mitigation}}}}\\[0.6mm]
\href{https://doi.org/10.1002/hsr2.72886}{Karamzad} & Ischaemic heart disease burden in MENA 1990-2021; PM attributable fraction 28.2\% & \textit{Health Sci Rep} & GBD secondary analysis & Ambient PM & MENA\\
\href{https://doi.org/10.1371/journal.pone.0347989}{Rees} & HEPA filtration to reduce respiratory infection in care homes & \textit{PLoS One} & RCT process evaluation & Filtered PM & UK\\
\href{https://doi.org/10.1016/j.jes.2026.01.009}{Shan} & Policy-driven emission reductions and post-pandemic rebound, high-resolution inventories & \textit{J Environ Sci (China)} & Emission inventory & PM emissions & China\\
\addlinespace[0.9mm]
\multicolumn{6}{@{}l@{}}{\cellcolor{Mist}\textsf{\textbf{\textcolor{Slate}{Other clinical endpoints}}}}\\[0.6mm]
\href{https://doi.org/10.3390/dj14070440}{Lupita} & PM$_{2.5}$, preventive-care capacity and caries/edentulism burden across the EU & \textit{Dent J} & Ecological panel & PM$_{2.5}$ & EU-27\\
\href{https://doi.org/10.1016/j.jdin.2026.06.008}{Park} & Long-term PM exposure and risk of androgenetic alopecia & \textit{JAAD Int} & Nationwide cohort & PM$_{2.5}$ / PM$_{10}$ & South Korea\\
\href{https://doi.org/10.3171/2026.3.PEDS2615}{Sun} & Environmental exposure and area deprivation vs. paediatric hydrocephalus outcomes & \textit{J Neurosurg Pediatr} & Registry cohort + geospatial & PM$_{2.5}$ (satellite-derived) & USA\\
\addlinespace[0.9mm]
\bottomrule
\end{longtable}}


\clearpage

% ==================================================================== PROVENANCE
\section*{\sffamily\bfseries\color{Deep}Provenance \& method}
\vspace{-1.5mm}{\color{Amber}\rule{\linewidth}{1.1pt}}\vspace{2.5mm}

\begin{keybox}
{\fontsize{7.2}{8.8}\selectfont
Entry window \textbf{2026-09-17 $\to$ 2026-09-17}, contiguous with the 16 September issue;
no gap. \textbf{This issue was built late}: the 17 September scheduled run did not fire, and
both 17 and 18 September were produced on 18 September as separate single-day issues rather
than merged into one two-day window, so the daily cadence and the per-day archive stay
intact.
\textbf{Sources.} PubMed \texttt{[Date - Entry]} run as two separate queries --- health axis
and instrumentation axis --- through both the local harvester and the connector:
\textbf{12 candidates, and the two clients agreed exactly} (10 health, 2 sensing, one record
on both legs), the first run in which neither client found anything the other missed.
Europe PMC \texttt{CREATION\_DATE} 31. Crossref windowed by ISSN across the tracked
non-PubMed journals \textbf{49 deposits}, which again carried every atmospheric-science
record in the issue --- the two \textit{ACP} papers, the \textit{AAQR} sensor framework and
the \textit{ES\&T} multi-model assessment are all invisible to PubMed. \textbf{arXiv
returned HTTP 406} on both windows, as on 15 and 16 September; per the 16 September finding
the leg is not repairable by header negotiation and remains dead. OpenAlex HTTP 429, as on
every recent run. No AGRICOLA records were returned, so the publication-year gate added on
16 September did not fire.
\textbf{Consensus returned nothing admissible, across two axes.} Sweep one (low-cost sensor
calibration, drift, co-location) and sweep two (satellite AOD--surface PM fusion, exposure
modelling), 20 records: \textbf{six already in the archive} on a title check, and
\textbf{all fourteen survivors failed an independent Crossref \texttt{created} check} ---
earliest \texttt{10.3390/s26010280} at 2026-01-02, latest \texttt{10.3390/math14132402} at
2026-07-06, with nothing inside ten weeks of the window. Four of the survivors
(\texttt{10.3390/atmos17040335}, \texttt{10.3390/s26030796},
\texttt{10.1109/access.2026.3656120}, \texttt{10.3390/s26010280}) were previously
\emph{rejected} in earlier issues rather than shipped, so they read as new against the
corpus but are not. Consensus has no entry-date filter; on this run its admitted yield was
zero for the second consecutive issue.
\textbf{Screening.} 92 raw records across the legs, 84 unique after cross-source dedupe, of
which 9 were caught by the PMID/DOI gate against \texttt{seen.json}.
\textbf{18 shipped, 57 excluded, 5 held.} Exclusions are overwhelmingly on-journal but
off-topic --- \textit{ES\&T} and \textit{Environmental Pollution} deposit heavily in soil,
water and PFAS chemistry, and \textit{Measurement Science and Technology} in machine vision
and inertial navigation, none of which touch airborne particles. The five held are same-day
deposits with no abstract recoverable from Crossref, Europe PMC or PubMed: three Elsevier
DOIs (\texttt{10.1016/j.atmosenv.2026.122363} WRF-Chem PM$_{2.5}$ forecasting,
\texttt{10.1016/j.scitotenv.2026.182313} Doha compound heat--pollution,
\texttt{10.1016/j.atmosres.2026.109340} Lecce aerosol source apportionment) and two Europe
PMC records with \textbf{no DOI deposited at all} (\textit{Frontiers in Toxicology} on
pollutant regulation of neurodegenerative markers, \textit{Environmental Science and
Ecotechnology} on tree equity and US urban health burden). All five are on-axis and are
deliberately kept out of \texttt{seen.json}; this is the third consecutive issue in which
the no-abstract-on-deposit pattern has cost tier-A candidates.
\textbf{Trial watch.} Registry sweep for the window is reported with the 18 September issue
and accumulated to \texttt{trials.json}, whose counter is derived from persisted NCT records
rather than incremented, per the 16 September reconciliation.
\textbf{Labels.} Subtopic, design, particle-metric and geography are assigned from title,
abstract, keywords and MeSH; single-label by dominant contribution, a judgement rather than
a controlled vocabulary. \textbf{Seventeen design strings were drafted and fourteen were
rewritten to reuse names the archive already owned} --- the failure the 15 September log
recorded as inventing fresh wording for existing designs, caught this time by diffing the
issue's design set against \texttt{design\_group} before the build rather than after. The
three genuinely new labels (adaptive mobile-robot sampling with CFD validation;
repeated-measures cohort with metabolite mediation; multi-city time-series with
random-effects meta-analysis) were registered in \texttt{plots.py} in the same edit that
wrote them onto records.
Effect estimates are transcribed verbatim from abstracts; the two Beijing percent changes
are converted to ratio equivalents and flagged in the forest-plot caption, and nothing else
is rescaled.
Content derived from PubMed and Crossref metadata; abstracts remain the copyright of their
respective publishers.\par}
\end{keybox}

\end{document}
